% GCQ — ODI 2026 Proposals-track submission (NeurIPS 2026 workshop)
% Compile: pdflatex main.tex (twice for references)
% NOTE: uses the official NeurIPS 2026 style file (neurips_2026.sty) placed
% next to this file. The \IfFileExists fallback below lets the draft compile
% without it, but the SUBMITTED pdf must use the official style.

\documentclass{article}

\PassOptionsToPackage{numbers}{natbib}
\IfFileExists{neurips_2026.sty}{%
  \usepackage{neurips_2026} % anonymized (default) — correct for double-blind submission
}{%
  % --- fallback so the draft compiles anywhere; replace with official style before submission ---
  \usepackage[margin=1in]{geometry}
  \newcommand{\fallbacknote}{\begin{center}\fbox{\parbox{0.9\linewidth}{\centering
    \textbf{DRAFT} — compiled without \texttt{neurips\_2026.sty}; page count is NOT meaningful in this mode.}}\end{center}}
}

\usepackage[utf8]{inputenc}
\usepackage[T1]{fontenc}
\usepackage{amsmath,amssymb}
\usepackage{booktabs}
\usepackage{graphicx}
\usepackage{xcolor}
\usepackage{hyperref}

\newcommand{\bbox}{\texttt{bbox\_2d}}

\title{GCQ: Measuring and Closing the Grounding Gap\\in Quantized Edge VLMs, Training-Free}

% Double-blind: author block intentionally empty (neurips_2026 default is anonymized).
\author{Anonymous}

\begin{document}

\maketitle
\IfFileExists{neurips_2026.sty}{}{\fallbacknote}

\begin{abstract}
Vision-language models (VLMs) reach phones, robots, and embedded systems almost exclusively through post-training quantization---yet the capability edge deployment usually exists for---precise spatial grounding---is the one nobody measures: every major VLM quantization method validates on VQA, OCR, and reasoning alone. We present the first systematic measurement of what quantization does to visual grounding, quantizing Qwen3-VL-2B across bit-widths, quantizers, and activation settings, and scoring grounding beside VQA and hallucination against image-blind floors. Three findings tell one story. \textbf{(1) The damage is real but hides:} at 4-bit, aggregate benchmarks degrade proportionally, but held-out detection stratified by object size reveals a fine-spatial-precision gradient---small-object recall $-53$\%, large-object unchanged---invisible to every standard benchmark; at GPTQ 3-bit, grounding dies outright (0\% boxes) while VQA partially survives. \textbf{(2) The damage is mappable:} a forward-only promote-one-module probe localizes grounding sensitivity in a mid-network attention band, uncorrelated with the VQA sensitivity map ($\rho=0.001$) and specific to the deployed quantizer ($\rho=0.08$). \textbf{(3) The damage is cheaply recoverable, training-free:} promoting the band to 8-bit ($+6$\% LLM weight memory) buys back half the lost grounding, beats random and VQA-driven allocation at identical memory, transfers to held-out detection, and leaves VQA unharmed. Because nothing is retrained, the recovery cannot silently trade away general capability; the recipe (GCQ) is drop-in for standard pipelines and doubles as an audit: compressed-VLM releases should report grounding on the model card.
\end{abstract}

\section{Introduction}

Every major VLM post-training-quantization (PTQ) paper---MBQ~\cite{mbq}, Q-VLM~\cite{qvlm}, VLMQ~\cite{vlmq}, LUQ~\cite{luq}, SPEED-Q~\cite{speedq}---validates on VQA, OCR, chart, and reasoning benchmarks (verified by full-text search of each), and SPEED-Q does so while explicitly targeting smartphone and robot deployment. With one incidental exception---GWQ~\cite{gwq} includes a single RefCOCO table to validate its general outlier-preserving method---none reports any localization metric or analyzes localization retention against VQA retention. Meanwhile, token pruning is known to collapse localization first (95\% VQA vs.\ 47\% grounding retention)~\cite{nuwa,feather}; quantization, the dominant edge-compression technique, is an unguarded blind spot. Since on-device deployment usually exists for \emph{perception}, practitioners may be shipping broken spatial understanding behind healthy VQA numbers on the model card: the default tags of Ollama's \texttt{qwen2.5vl} (4.1M pulls) serve 4-bit checkpoints on a page that advertises bounding-box localization~\cite{ollamaqwen}.

That is a threefold gap --- no systematic \emph{measurement} of grounding under quantization, no \emph{method} that allocates precision by capability, and no \emph{reporting practice} that would expose the damage --- and this proposal answers each part in turn:
\begin{enumerate}
  \item \textbf{A measurement study (the centerpiece).} Quantize Qwen3-VL-2B-Instruct across bit-widths and measure the \emph{grounding gap}: grounding retention vs.\ VQA retention, side by side, for the first time to our knowledge. (\emph{Grounding} = referring-expression localization; held-out ODinW probes transfer to detection.)
  \item \textbf{One training-free intervention.} \emph{Localization-sensitivity-guided mixed precision}: profile how much each module's quantization damages grounding, then promote the most grounding-sensitive modules to 8-bit under a small memory budget---forward-only, greedy, over a completely standard quantization pipeline.
  \item \textbf{An evaluation-protocol claim.} Selection on grounding metrics with VQA as a constraint---inverting standard practice---and grounding reported on compressed-model cards.
\end{enumerate}

We call the recipe \textbf{GCQ} (grounding-consistent quantization): one model, one intervention, one protocol, no training anywhere---a fine-tune can restore grounding by quietly trading away VQA (LoRA alone hurts grounding~\cite{miniinternvl}); a training-free pipeline cannot.

\section{Background and related work}

\paragraph{The asymmetry is documented---for pruning.} FEATHER~\cite{feather} showed that visual-token pruning leaves most benchmarks intact while localization accuracy decays roughly linearly to zero, and that popular benchmarks fail to detect this because they require minimal grounding. N\"uwa~\cite{nuwa} quantified the asymmetry (the 95\%-vs-47\% figure above, under 88.9\% token reduction); further pruning studies report the same fragility for pixel and UI grounding~\cite{litelvlm,focusui}. The mechanism is intuitive: a correct box needs a preserved spatial reference frame and fine positional detail; ``what color is the dog'' does not.

\paragraph{For quantization, nobody has looked systematically.} The VLM PTQ literature evaluates exclusively on VQA/OCR/chart/doc/reasoning; MBQ~\cite{mbq} frames sensitivity as a \emph{modality} issue, not a \emph{capability} issue---no existing method asks which layers carry \emph{localization}. Direct evidence is two isolated datapoints pointing in opposite directions: GWQ's incidental near-lossless RefCOCO table at ${\sim}4.6$ bits with FP16 outliers retained~\cite{gwq}---whose own \emph{baseline} rows show 4-bit AWQ dropping RefCOCO by 4.9 points, unremarked---vs.\ an applied study's 4-bit Qwen-VL that classifies damage correctly but ``fail[s] consistently'' to localize it~\cite{vehicledamage}. The Qwen3 family is known to degrade steeply at $\le$3 bits~\cite{qwen3quant}; component-wise vision/LLM bit allocation is evaluated on MME only~\cite{rethinkingsmall}, SQAP-VLA~\cite{sqapvla} protects visual \emph{input} tokens, and MODE~\cite{mode} allocates expert bits by modality. Closest to this work, MABA~\cite{maba} performs gradient-guided VLM mixed-precision allocation (LLaVA/Qwen-VL/InternVL)---so GCQ is not the first VLM mixed-precision or gradient-guided allocator. MABA allocates by aggregate \emph{modality} reconstruction error and reports only general multi-modal benchmarks, with no grounding, localization, or detection evaluation. A second prior~\cite{taq} allocates per-layer precision by task, but for text-only LLMs. GCQ's novelty is narrower: allocation conditioned on one \emph{capability} (grounding), a coordinate-token sensitivity signal, and object-scale/held-out detection evaluation---none present in either prior.

\paragraph{Why Qwen3-VL-2B.} Qwen3-VL~\cite{qwen3vl} trains grounding natively and emits coordinates as plain text in a normalized $[0,1000]$ system, and the 2B is edge-realistic (${\sim}1.5$\,GB at W4; we write WxAy for $x$-bit weights, $y$-bit activations). Following~\cite{rethinkingsmall}, the vision tower stays at BF16/INT8; all aggressive quantization targets the LLM.

\section{Measuring the grounding gap}

\paragraph{Pipeline overview.} Freeze hygienic evaluation subsets; quantize the LLM (RTN/GPTQ, W8--W3, W4A8 stress); score grounding, VQA, and hallucination from \emph{decoded} outputs against image-blind floors; profile per-module sensitivity with promote-one probes; allocate extra precision greedily under a memory budget; verify against matched-memory controls, a held-out OOD set, and full test splits.

\paragraph{Setup and data hygiene.} Reference model $=$ Qwen3-VL-2B-Instruct at BF16; quantized model $=$ the same with quantized LLM weights. $D_{\text{probe}}$: ${\sim}512$ RefCOCO-train samples in the chat template, used only for sensitivity profiling (\S\ref{sec:method}). $D_{\text{dev}}$: 1{,}000 RefCOCO-val samples, used \emph{only} for selection---never for calibration or profiling. Reported RefCOCO-val numbers exclude $D_{\text{dev}}$; ODinW-13 is never touched by calibration, profiling, or selection---probe, selection, and test data are strictly disjoint. RefCOCOg uses the leak-free umd split, and $D_{\text{probe}}$ excludes every image appearing in any RefCOCO variant's val/test splits (the variants draw on overlapping COCO images).

\paragraph{Evaluation pipeline.} Each REC sample is one image plus one chat-template query---``Locate the \{expression\}, output its \bbox{} in JSON''---decoded greedily, parsed (regex fallback), mapped $[0,1000] \to$ pixels, and scored against the ground-truth box: correct at IoU $\ge 0.5$; mean GIoU recorded. Unparseable outputs count as wrong; the parse-failure rate is reported separately (quantization can break the output format itself). For ODinW-13, one category per prompt (``locate every \{class\}''), greedy IoU matching, precision/recall@0.5. Detection recall is additionally stratified by ground-truth box size, where degradation is expected to concentrate.

\paragraph{Detection in the loop.} Every choice a compression pipeline normally makes by perplexity or VQA---group size, bit budget, which modules get promoted---is made here by $D_{\text{dev}}$ grounding metrics (REC acc@0.5 $+$ mean GIoU). General capability is a \emph{constraint}: VQA and hallucination scores must stay within a pre-registered tolerance (1.5 points) of the same-method text-calibrated quantized baseline, checked on the fixed VQAv2 and POPE subsets.

\section{GCQ: localization-sensitivity-guided mixed precision}
\label{sec:method}

\paragraph{Step 1---profiling (forward-only).} Profiling runs in the same direction the allocator spends: start from the \emph{fully quantized} model at the aggressive width (a standard text-calibrated checkpoint---the base PTQ pipeline is deliberately untouched), then for each LLM module $m$ (attention and MLP per block; ${\sim}2{\times}28 \approx 56$ modules on the 2B) promote \emph{only} $m$ to 8-bit and score the \emph{reduction} in grounding damage on $D_{\text{probe}}$ with a decoding-free proxy:
\begin{equation}
s_m \;=\; \mathrm{KL}_{\text{base}} - \mathrm{KL}_{+m}, \qquad
\mathrm{KL}_{c} \;=\; \mathbb{E}_{D_{\text{probe}}}\Big[\textstyle\sum_{t \in C} \mathrm{KL}\big(p_{\text{ref}}(\cdot \mid x, y_{<t}) \,\|\, p_{c}(\cdot \mid x, y_{<t})\big)\Big],
\end{equation}
where $C$ is the set of coordinate-token positions in the teacher-forced target (the tokens covering digits, commas, and brackets inside \bbox{} arrays; well-defined under multi-digit tokenizer merges). This measures exactly the marginal utility the greedy allocator consumes, at the deployment operating point. Because token-level agreement correlates imperfectly with geometric error~\cite{groundingtokens}, the proxy is never trusted on its own: a decoded 9-module spot-check is directionally consistent but underpowered (single-module effects sit inside sampling noise), so the operative validation is allocation-level---no configuration is accepted on proxy numbers (Step 2).

\paragraph{Step 2---allocation (greedy, two-tier).} Given an average-bits budget $B$, promote the module with the highest $s_m$ per extra byte to 8-bit until the budget is exhausted; embeddings and \texttt{lm\_head} stay 8-bit always (they carry the digit vocabulary). The arithmetic, stated honestly: at $B{=}4.25$ the budget promotes ${\sim}88$M of the ${\sim}1.4$B quantizable parameters (${\sim}{+}6\%$ LLM weight memory, ${\sim}$50--60\,MB, over uniform W4); $B{=}4.5$ doubles both; comparisons are always between allocation \emph{policies} at matched memory. Known proxy--deployment gaps (per-module isolation, teacher forcing, RTN-for-profiling vs.\ GPTQ-final) are all bounded by one rule: no allocation is accepted on proxy numbers---every candidate is evaluated end-to-end on decoded $D_{\text{dev}}$ metrics; the proxy only orders the search. Deliberately absent: calibration-set changes (gains attribute to allocation alone), training of any kind, learned allocators, architecture or tokenizer changes.

\section{Results}
\label{sec:results}

\paragraph{Settings.} RTN and our validated GPTQ implementation, weight-only, groupwise symmetric $g{=}128$, LLM blocks only; W8/W4/W3; W4A8 adds simulated per-token dynamic 8-bit activations; GCQ at $B \in \{4.25, 4.5\}$. All rows share one evaluation pipeline on frozen seed-0 subsets (REC 1k, VQAv2 5k, POPE 9k, ODinW 500 images). Retention is floor-corrected, $R' = (S - S_{\text{floor}})/(S_{\text{BF16}} - S_{\text{floor}})$, against image-blind floors (REC 3.0, VQAv2 36.8, POPE 50.0); tables show absolute scores.

\begin{table}[t]
  \centering
  \caption{\textbf{The grounding gap across quantization settings.} Absolute scores; identical pipeline every row; ODinW-13 is held out from all calibration, profiling, and selection; W3 rows are analyzed in the text.}
  \label{tab:gap}
  \small
  \begin{tabular}{llcccc}
    \toprule
    Setting & Avg bits & REC@0.5 & VQAv2 & POPE & ODinW F1 \\
    \midrule
    BF16 reference & 16 & 88.9 & 81.1 & 89.4 & 0.654 \\
    RTN W8 & 8 & 88.7 & 81.3 & 89.4 & 0.661 \\
    \midrule
    RTN W4 & 4 & 84.5 & 78.0 & 88.8 & 0.582 \\
    GPTQ W4 & 4 & 85.4 & 78.1 & 89.4 & 0.526 \\
    RTN W4A8 & 4 & 82.8 & 77.5 & 88.7 & 0.573 \\
    \midrule
    RTN W3 & 3 & 0.0 & 1.3 & 50.4 & 0.000 \\
    GPTQ W3 & 3 & 0.0 & 17.7 & 59.4 & 0.000 \\
    \bottomrule
  \end{tabular}
\end{table}

\begin{table}[t]
  \centering
  \caption{\textbf{Allocation policy at matched memory} (avg $4.25$ bits $=$ uniform W4 $+$ ${\sim}88$M parameters promoted to 8-bit, $+6$\% LLM weight memory; $B{=}4.5$ doubles the premium). Bottom block: the same recipe on GPTQ with a GPTQ-profiled map.}
  \label{tab:alloc}
  \small
  \begin{tabular}{lccc}
    \toprule
    Policy (base quantizer RTN) & REC@0.5 & mean GIoU & VQAv2 \\
    \midrule
    Uniform W4 (no promotion) & 84.5 & 0.754 & 78.0 \\
    Random promotion (3 seeds) & 84.9$\pm$0.3 & 0.758 & 78.2$\pm$0.4 \\
    VQA-sensitivity-driven & 85.1 & 0.764 & 78.6 \\
    \textbf{GCQ (ours)} & \textbf{86.6} & \textbf{0.775} & \textbf{78.4} \\
    \textbf{GCQ at $B{=}4.5$ (ours)} & \textbf{87.3} & \textbf{0.781} & \textbf{78.7} \\
    \midrule
    Uniform GPTQ-W4 & 85.4 & 0.745 & 78.1 \\
    \textbf{GCQ on GPTQ (matched map)} & \textbf{85.5} & \textbf{0.768} & \textbf{78.2} \\
    \bottomrule
  \end{tabular}
\end{table}

\begin{figure}[t]
  \centering
  \includegraphics[width=\linewidth]{fig1_retention.pdf}
  \caption{\textbf{Measured retention and allocation (RTN, Qwen3-VL-2B).} (a) Floor-corrected retention: capabilities degrade proportionally at W4 and collapse together at W3 --- no pruning-style grounding-first asymmetry in the aggregate; GCQ (diamonds) recovers grounding at small extra memory. (b) At an identical $+88$M-parameter budget, grounding-driven allocation beats random (${\sim}5.7{\times}$ the recovery) and VQA-driven (${\sim}3.5{\times}$) promotion.}
  \label{fig:main}
\end{figure}

\begin{figure}[t]
  \centering
  \includegraphics[width=\linewidth]{fig2_sensitivity.pdf}
  \caption{\textbf{Capability sensitivity maps are independent} (Spearman $\rho = 0.001$, top-7 overlap $0/7$). Grounding sensitivity (top) concentrates in a mid-network attention band (layers 10--17); VQA sensitivity (bottom) concentrates early; the most grounding-critical module (layer-15 attention) is \emph{negative} for VQA.}
  \label{fig:sens}
\end{figure}

\paragraph{The gap is real, hidden by aggregate benchmarks, and fatal at 3 bits.} At W4 the aggregate metrics degrade roughly proportionally---floor-corrected retention POPE 98.4\% $>$ grounding 94.9\% $>$ VQA 92.8\% (Table~\ref{tab:gap}, Figure~\ref{fig:main}a)---no pruning-style asymmetry. But the aggregate hides a gradient: on held-out ODinW stratified by ground-truth box size, W4 costs $-53$\% small-object recall ($18.5 \to 8.7$), $-26$\% medium, and large objects are untouched. The asymmetry was never absent; it is \emph{invisible to large-object benchmarks}---and to VQA and POPE entirely. Out of domain the damage is bigger: 11\% of relative ODinW F1 lost vs.\ 5\% of in-domain REC. At 3 bits the capabilities separate outright: RTN-W3 collapses globally (every score at its floor), but GPTQ-W3 half-survives, and there \emph{grounding dies first}---0\% parseable boxes in and out of domain, while VQA and POPE retain partial function (17.7 and 59.4, the latter above its 50.0 floor; answers verified varied, not a yes-bias). The collapse is not band-reversible: the grounding band at 8-bit over W3 more than doubles VQA ($17.7 \to 45.1$) yet yields zero boxes---grounding is not merely first to die, it is the costliest to save. W8 is lossless on every metric; W4A8 costs a further $-1.7$ REC over weight-only W4. 
\paragraph{The damage is mappable---and the map is capability- and quantizer-specific.} The promote-one-module profile over all 56 module-groups (Figure~\ref{fig:sens}) is concentrated, not diffuse: it peaks in a contiguous band of \emph{attention} modules in layers 10--17 of 28, and MLPs lose every per-byte comparison at the $B{=}4.25$ budget---converging with mechanistic evidence that localization concentrates in few attention heads~\cite{locheads,causalheads}. The VQA map is statistically unrelated (Spearman $\rho = 0.001$; top-7 module overlap $0/7$), so protecting grounding and protecting VQA are \emph{different} interventions. The map is also quantizer-specific ($\rho$ between RTN- and GPTQ-profiled maps $= 0.082$): transplanting the RTN map onto GPTQ \emph{costs} $-1.2$ REC with 2.8\% parse failures, while the matched GPTQ-profiled map restores parity, recovers 44\% of the lost GIoU with zero parse failures (Table~\ref{tab:alloc}, bottom), and carries out of domain (ODinW F1 $0.526 \to 0.595$, 54\% of GPTQ's held-out loss). The prescription: profile on the quantizer you deploy.

\paragraph{Well-placed bits buy grounding back---training-free.} At identical memory (Table~\ref{tab:alloc}), GCQ recovers $+2.1$ REC over uniform W4---48\% of the W4 grounding loss---versus $+0.4\pm0.3$ for random promotion and $+0.6$ for VQA-driven promotion; $B{=}4.5$ recovers 64\%, monotonically. The constraint holds with margin: GCQ's VQAv2 is $+0.4$ ($B{=}4.25$) and $+0.7$ ($B{=}4.5$) \emph{above} the uniform baseline, and its POPE (88.8) matches uniform W4 exactly. The RefCOCO-chosen band transfers out of domain: ODinW F1 $0.582 \to 0.626$ at $B{=}4.25$ and $0.633$ at $B{=}4.5$ (61\% and 71\% of the held-out loss), with small/medium-object recall recovery of 55--58\%. Everything replicates on the full official test splits ($n{=}5657/5095$): 49--60\% recovery on testA, 32--43\% on the harder testB, significant under a paired bootstrap (testA $+1.84$, 95\% CI $[1.29, 2.39]$).

\paragraph{Consequences.} Practitioners: a stock 4-bit VLM has silently lost half its small-object detection---audit before shipping; $+6\%$ well-placed memory buys half back at zero VQA cost. Toolchains: profile sensitivity on the deployed quantizer; report grounding on model cards. Research: capabilities are separable, mappable structures---compression can become capability-aware, and the ${\sim}35$--$50\%$ residual is the target for recovery training.

\section{Hypotheses, controls, and outcomes}

\paragraph{Pre-registered hypotheses, resolved.} \textbf{H1} (grounding trails VQA by $>$10 floor-corrected points at W4): \emph{refuted in the aggregate}---the gap is a size gradient plus a W3 grounding-first collapse, sharper claims than H1 anticipated. \textbf{H2} (sensitivity is structured and differs from VQA's): \emph{confirmed} ($\rho=0.001$, $0/7$ overlap). \textbf{H3} ($\ge$90\% recovery at $B\le4.5$): \emph{not reached}---48--64\%; the pre-registered escalation (distillation, deferred) covers the remainder. Reporting refuted and partial outcomes is deliberate.

\paragraph{Controls.} Matched memory for every allocation comparison; same calibration data and pipeline for every quantizer row; selection only on $D_{\text{dev}}$, disjoint from all reported numbers; held-out ODinW touched once; image-blind floors; parse failures reported separately.

\paragraph{Known limitations.} (i) $s_m$ is teacher-forced and per-module; proxy validity rests on the end-to-end decoded selection and allocation-level controls. (ii) The novelty claim is \emph{to our knowledge} (closest: \cite{gwq,mbq,feather,nuwa,sqapvla,rethinkingsmall,maba}); GCQ does not claim to be the first VLM mixed-precision or gradient-guided allocator (that is MABA~\cite{maba}), only the first conditioned on a single capability with coordinate-specific sensitivity and grounding evaluation. (iii) Single-model study; the 4B/8B scale study is deferred. The campaign cost ${\sim}$50 GPU-hours; all artifacts will be released.

\begin{thebibliography}{20}\small
\bibitem{qwen3vl} Bai et al. Qwen3-VL technical report. arXiv:2511.21631.
\bibitem{feather} Endo et al. Feather the throttle: Revisiting visual token pruning for vision-language model acceleration. arXiv:2412.13180, ICCV 2025.
\bibitem{nuwa} Huang et al. N\"uwa: Mending the spatial integrity torn by VLM token pruning. arXiv:2602.02951.
\bibitem{litelvlm} Lee et al. CLIP tricks you: Training-free token pruning for efficient pixel grounding in large vision-language models. arXiv:2605.13178.
\bibitem{focusui} Ouyang et al. FocusUI: Efficient UI grounding via position-preserving visual token selection. arXiv:2601.03928.
\bibitem{mbq} Li et al. MBQ: Modality-balanced quantization for large vision-language models. arXiv:2412.19509, CVPR 2025.
\bibitem{qvlm} Wang et al. Q-VLM: Post-training quantization for large vision-language models. arXiv:2410.08119, NeurIPS 2024.
\bibitem{vlmq} Xue et al. VLMQ: Token saliency-driven post-training quantization for vision-language models. arXiv:2508.03351.
\bibitem{luq} Bhatnagar et al. LUQ: Layerwise ultra-low bit quantization for multimodal large language models. arXiv:2509.23729.
\bibitem{speedq} Guo et al. SPEED-Q: Staged processing with enhanced distillation towards efficient low-bit on-device VLM quantization. arXiv:2511.08914.
\bibitem{gwq} Shao et al. GWQ: Gradient-aware weight quantization for large language models. arXiv:2411.00850.
\bibitem{qwen3quant} Zheng et al. An empirical study of Qwen3 quantization. arXiv:2505.02214.
\bibitem{rethinkingsmall} Shin et al. Rethinking small VLM quantization: From component-wise analysis to hardware-aware edge deployment. arXiv:2607.08029.
\bibitem{locdistill} Zheng et al. Localization distillation for dense object detection. arXiv:2102.12252.
\bibitem{sqapvla} Fang et al. SQAP-VLA: A synergistic quantization-aware pruning framework for vision-language-action models. arXiv:2509.09090.
\bibitem{groundingtokens} Ren et al. Grounding everything in tokens for multimodal large language models. arXiv:2512.10554.
\bibitem{vehicledamage} Hogale et al. Grounding agentic VLMs with dedicated segmentation for fine-grained vehicle damage assessment. arXiv:2608.02470.
\bibitem{miniinternvl} Gao et al. Mini-InternVL: A flexible-transfer pocket multimodal model with 5\% parameters and 90\% performance. arXiv:2410.16261.
\bibitem{taq} LeVi et al. You had one job: Per-task quantization using LLMs' hidden representations. arXiv:2511.06516.
\bibitem{mode} Chen et al. MODE: Modality-decomposed expert-level mixed-precision quantization for MoE multimodal LLMs. arXiv:2606.17118.
\bibitem{maba} Zhang, Zhu, Wang, Wu, Lin. Modality-aware bit allocation for mixed-precision quantization of vision-language models. CVPR 2026 Findings, pp.\ 9305--9315.
\bibitem{locheads} Kang et al. Your large vision-language model only needs a few attention heads for visual grounding. arXiv:2503.06287, CVPR 2025.
\bibitem{causalheads} Schauml\"offel et al. Mechanisms of object localization in vision-language models. arXiv:2605.19792, CVPR 2026.
\bibitem{ollamaqwen} Ollama model library: qwen2.5vl. \url{https://ollama.com/library/qwen2.5vl}, accessed Aug 2026.
\end{thebibliography}

\subsubsection*{LLM usage disclosure}
A large language model was used as a writing and editing aid in preparing this proposal; all technical content, claims, and experimental design are the authors' own, and the authors take full responsibility for the submission's content.

\appendix
\section{Deferred to the full paper}
Grounding-calibrated PTQ (with a three-arm attribution ablation separating grounding-specific calibration effects from generic multimodal-calibration effects); the coordinate-weighted distillation recovery~\cite{locdistill} for the residual unrecovered grounding; the 4B/8B scale study; prompted-detection COCO AP with size-stratified reporting; finer bit-budget Pareto sweeps; on-device latency/memory export measurements. A small-object-weighted probe (smallest-quartile boxes) selects a 5/7-overlapping allocation but does not improve small-object recovery at matched memory, marking the current training-free floor.

\end{document}
